\documentclass[10pt,a4paper]{article}
\usepackage[T1]{fontenc}
\usepackage[utf8]{inputenc}
\usepackage[margin=2.2cm]{geometry}
\usepackage{booktabs}
\usepackage{longtable}
\usepackage{array}
\usepackage{xcolor}
\usepackage{amsmath}
\usepackage[hidelinks]{hyperref}

\definecolor{hdr}{HTML}{1F3864}
\makeatletter
\renewcommand\section{\@startsection{section}{1}{\z@}%
  {-2.4ex \@plus -1ex \@minus -.2ex}{1.3ex \@plus.2ex}%
  {\normalfont\large\bfseries\color{hdr}}}
\renewcommand\subsection{\@startsection{subsection}{2}{\z@}%
  {-2.0ex \@plus -1ex \@minus -.2ex}{1.0ex \@plus.2ex}%
  {\normalfont\normalsize\bfseries\color{hdr}}}
\makeatother
\newenvironment{tightitem}{%
  \begin{itemize}\setlength{\itemsep}{1pt}\setlength{\parskip}{0pt}%
  \setlength{\topsep}{3pt}\setlength{\parsep}{0pt}}{\end{itemize}}
\renewcommand{\arraystretch}{1.15}

\title{\vspace{-1.4cm}\color{hdr}\bfseries Estimating Relatedness in Admixed and Structured Populations\\[2pt]
\large\mdseries\color{black} An overview of existing methods}
\author{}
\date{\small Compiled \today}

\begin{document}
\maketitle
\vspace{-1.2cm}

\section{The problem}

Classical relatedness estimators --- the method-of-moments estimator in PLINK, or maximum-likelihood
estimation of the Cotterman coefficients $(k_0,k_1,k_2)$ --- assume that both individuals are drawn from a
single, homogeneous population with known allele frequencies $p$. Under admixture this assumption fails in
two ways:

\begin{tightitem}
  \item \textbf{Population-level allele frequencies are the wrong frequencies.} Two unrelated individuals
  who share ancestry from the same source population share alleles more often than the pooled frequency
  predicts. They are scored as related; kinship is inflated, and the inflation scales with how differentiated
  the source populations are ($F_{ST}$).
  \item \textbf{Frequencies are individual-specific.} For an individual with ancestry proportions
  $q_{i}=(q_{i1},\dots,q_{iK})$ and ancestry-specific frequencies $f_{1},\dots,f_{K}$, the relevant
  frequency is $\pi_{i}=\sum_{k} q_{ik} f_{k}$, which differs between the two members of a pair. A pair
  with \emph{different} ancestry profiles is then scored as \emph{less} related than unrelated --- KING-robust,
  for instance, returns kinship estimates around $-0.13$ for unrelated pairs of dissimilar ancestry.
\end{tightitem}

Essentially all methods below are answers to the same question: \emph{what do you condition on so that
``sharing alleles'' means ``sharing recent ancestors'' rather than ``sharing continental ancestry''?}
Three families of answers exist, and they are summarised in Table~\ref{tab:methods}.

\section{Three strategies}

\subsection{(A) Condition on estimated ancestry: ancestry-specific allele frequencies}

Run a global ancestry / admixture model first (ADMIXTURE, STRUCTURE, frappe, NGSadmix) to obtain
$\hat{q}_{i}$ and $\hat{f}_{k}$, then plug the individual-specific frequencies
$\hat{\pi}_{i}=\sum_k \hat q_{ik}\hat f_{k}$ into a relatedness model. This is the most explicit approach and
gives interpretable IBD-state output, at the cost of depending on $K$ being sensible and on the admixture
model fitting.

\begin{tightitem}
  \item \textbf{REAP} (Thornton et al.\ 2012) --- moment estimator of $(k_0,k_1,k_2)$ and kinship using
  individual-specific frequencies conditional on genome-wide ancestry. The first method to make this idea
  standard; validated under ancestry-related assortative and disassortative mating.
  \item \textbf{RelateAdmix} (Moltke \& Albrechtsen 2014) --- full maximum-likelihood estimation of
  $(k_0,k_1,k_2)$ given $\hat q$ and $\hat f$, accounting for the fact that the two IBD alleles must come
  from the \emph{same} source population. R package plus a multithreaded C++ command-line tool.
  \item \textbf{InRelate} (Sethuraman 2018) --- ML over the full nine-state Jacquard IBD framework
  (so inbred as well as outbred pairs), for multi-allelic markers, given admixture-model ancestry.
  \item \textbf{SEEKIN} (Dou et al.\ 2017) --- kinship for shallow / off-target sequencing and array data;
  models genotype uncertainty (dosages from imputation) and corrects for structure using ancestry
  coordinates. Works down to $\sim$0.15$\times$ depth; \texttt{SEEKIN-het} is the admixture-aware mode.
  \item \textbf{NGSremix} (N\o hr et al.\ 2021) --- ML estimation of $(k_0,k_1,k_2)$ for admixed pairs from
  \emph{genotype likelihoods}, so genotypes are never called. Designed for low-depth NGS ($\leq$4$\times$),
  where calling genotypes would itself bias relatedness; outperforms PLINK, KING, RelateAdmix and NgsRelate
  in that regime.
\end{tightitem}

\subsection{(B) Condition on principal components: PC-derived individual-specific frequencies}

Instead of a discrete $K$-population model, regress genotypes on ancestry-representative PCs to predict
each individual's expected allele frequency, and define relatedness from the residuals. No choice of $K$,
no assumption that ancestral populations are discrete --- this is the ``model-free'' route.

\begin{tightitem}
  \item \textbf{PC-AiR} (Conomos, Miller \& Thornton 2015) --- \emph{not} a relatedness estimator, but the
  necessary first half. Ordinary PCA in a sample containing relatives captures family structure rather than
  ancestry. PC-AiR partitions the sample into an ancestry-representative unrelated subset, runs PCA there,
  and projects the relatives back, giving PCs that reflect ancestry and not pedigree.
  \item \textbf{PC-Relate} (Conomos, Reiner, Weir \& Thornton 2016) --- the second half. Uses the PC-AiR PCs
  to obtain individual-specific allele frequencies, then estimates kinship coefficients, IBD-sharing
  probabilities and inbreeding coefficients from the residual genotype correlations.
  \item In practice the two are run as one iterative pipeline (PC-AiR $\rightarrow$ PC-Relate $\rightarrow$
  refined unrelated set $\rightarrow$ PC-AiR again), implemented in the Bioconductor package
  \textbf{GENESIS}. This is the de-facto standard in large human cohorts (TOPMed, UK Biobank-style analyses).
\end{tightitem}

\subsection{Where the individual allele frequencies come from: PCAngsd and EMU}

Both strategy (A) and strategy (B) reduce to the same quantity: the \emph{individual allele frequency}
$\pi_{is}$, the expected frequency at site $s$ for individual $i$ given their ancestry. Strategy (A) gets it
from a discrete $K$-population decomposition, $\pi_{is}=\sum_k q_{ik}f_{ks}$; strategy (B) gets it from a
low-rank (PCA) decomposition of the genotype matrix. Two tools from the low-rank side deserve separate
mention, because they are what make the whole framework usable on data where genotypes cannot be trusted.

\textbf{PCAngsd} (Meisner \& Albrechtsen 2018) estimates individual allele frequencies for \emph{low-depth}
NGS data directly from genotype likelihoods, so no genotype is ever called. The model is
$\Pi = \mathbf{S}\mathbf{A}$: individual frequencies are a low-rank product of a population-structure
matrix and a loading matrix. It is fitted by an iterative EM-like scheme:
\begin{enumerate}
  \item estimate population frequencies $\hat p_s$ from the genotype likelihoods;
  \item form posterior genotype probabilities $P(G_{is}=g\mid X_{is},\hat\pi_{is})$ by Bayes' rule and take
  \emph{dosages} $E[G_{is}\mid X_{is},\hat\pi_{is}]$;
  \item centre the dosages, truncate their SVD at $D$ components, and add $2\hat p_s$ back to obtain
  $\hat\Pi$;
  \item feed $\hat\pi_{is}$ back in as the prior for step 2, and iterate to convergence.
\end{enumerate}
The covariance matrix is then computed from dosages centred on $2\hat p_s$, which is valid under structure
precisely \emph{because} the dosages already condition on $\hat\pi$. The same $\hat\Pi$ is then reused
downstream: admixture proportions by NMF, $\hat\Pi \approx QF^{\!\top}$ (\texttt{-{}-admix}), which is exactly
the $\hat q,\hat f$ that RelateAdmix and NGSremix need, plus inbreeding coefficients, HWE tests, selection
scans and genotype calling. In other words, PCAngsd supplies the input for strategy~(A) as well as the PCs of
strategy~(B) from a single fit.

One caveat, since it is natural to assume otherwise: \textbf{PCAngsd does not itself estimate pairwise
kinship any more.} Versions up to 0.99 had a \texttt{-kinship} flag implementing a PC-Relate-style (Conomos)
kinship matrix adapted to genotype likelihoods, plus a \texttt{-relate} option for filtering related
individuals out of a Beagle file. Both were dropped in the January 2021 Cython rework, and the current
release (v1.36) exposes only \emph{inbreeding} estimators (\texttt{-{}-inbreed-sites},
\texttt{-{}-inbreed-samples}). PCAone is the same: \texttt{-{}-inbreed}, but no kinship. In a present-day
pipeline PCAngsd and PCAone are therefore the \emph{front-end} producing $\hat\pi$ or $\hat q,\hat f$; the
relatedness estimate itself comes from NGSremix, RelateAdmix or PC-Relate.

\textbf{EMU} (Meisner, Liu, Huang \& Albrechtsen 2021) solves the neighbouring problem: PCA when the data are
\emph{called} genotypes but enormous fractions are \emph{missing} --- ultra-low-coverage sequencing, NIPT
data, merged arrays with different marker sets, ancient DNA. Standard PCA has to impute missing entries
(usually with the mean), and individuals then separate along PCs according to \emph{how much data they are
missing} rather than their ancestry. EMU uses an accelerated EM algorithm for PCA that models the
missingness explicitly and iteratively, alternating between a truncated SVD and re-imputation of the missing
entries from the current low-rank fit --- structurally the same idea as PCAngsd's loop, with missingness in
place of genotype uncertainty. It scales to biobank-sized data (Python implementation).

\textbf{PCAone} (Li, Meisner \& Albrechtsen 2023) is the scalable implementation of this whole layer. It is a
C++ PCA framework offering three fast SVD algorithms --- implicitly restarted Arnoldi (\texttt{-{}-svd 0}),
single-pass randomized SVD with power iterations (\texttt{-{}-svd 1}) and the proposed window-based
randomized SVD, winSVD (\texttt{-{}-svd 2}) --- each with in-core \emph{and} out-of-core modes, so biobank-scale
matrices need not fit in RAM. Crucially for the present purpose it re-implements \emph{both} the EMU and the
PCAngsd algorithms, so missingness and genotype-likelihood uncertainty are handled inside the same fast
framework, reading PLINK BED, PLINK2 PGEN, BGEN, Beagle and CSV. It also provides ancestry-adjusted LD
statistics for pruning and clumping --- relevant here because LD pruning with \emph{unadjusted} LD in a
structured sample preferentially removes exactly the ancestry-informative sites the relatedness estimator
needs to be conditioned on.

For relatedness work the practical point is the choice of front-end:
\begin{tightitem}
  \item \emph{low depth, genotype likelihoods available} $\rightarrow$ PCAngsd (or NGSadmix for discrete $K$),
  then NGSremix for the actual relatedness estimation;
  \item \emph{called genotypes, heavy and non-random missingness} $\rightarrow$ EMU for the ancestry axes,
  then a PC-based relatedness estimator;
  \item \emph{called genotypes, complete data} $\rightarrow$ PC-AiR, or ADMIXTURE for discrete $K$.
\end{tightitem}
A relatedness estimate is only as good as the $\hat\pi$ it conditions on, so this step is not a preliminary
--- it is half the method.

\subsection{Correlation of residuals as a kinship estimator: evalAdmix}

There is a fourth route into the same quantity, arrived at from the opposite direction. \textbf{evalAdmix}
(Garcia-Erill \& Albrechtsen 2020) was built to \emph{test the fit} of an admixture analysis, not to estimate
relatedness. Given genotypes and the output of ADMIXTURE/STRUCTURE/NGSadmix, it forms the residual
$g_{is}-2\hat\pi_{is}$ between the observed genotype and the genotype predicted by the admixture model, and
reports the matrix of pairwise \emph{correlations of residuals}. Under a correctly specified model the
residuals of two individuals are uncorrelated, so entries near zero indicate good fit; systematic positive
correlation within groups and negative correlation between groups indicates the model is wrong (too small a
$K$, ghost populations, drift after admixture, non-discrete ancestral populations).

But note what the residual is: exactly the deviation from the individual allele frequency. Two individuals
who share recent ancestors deviate from $\hat\pi$ in the \emph{same direction}, so relatedness produces a
positive residual correlation --- which is precisely PC-Relate's estimator, with the admixture model in place
of the PC regression. The correlation of residuals is therefore a kinship estimate; its documentation
flags this the other way round, warning that a positive correlation ``might also be due to relatedness''.
Concretely, the residual correlation estimates $r_{ij}=2\phi_{ij}$, the coancestry/relatedness coefficient,
not $\phi_{ij}$ itself.

This is easy to confirm. Running evalAdmix on the 126-individual example data distributed with RelateAdmix
($K=2$, 104{,}290 SNPs) and comparing against RelateAdmix's own ML estimates on the same data and same
$\hat q,\hat f$:

\begin{center}
\small
\begin{tabular}{@{}lrrr@{}}
\toprule
Relationship class & $\hat\phi$ (RelateAdmix) & mean corr.\ of residuals & $n$ pairs \\
\midrule
unrelated        & 0.000 & $-0.003$ & 7816 \\
second cousins   & 0.012 & 0.036 & 9 \\
first cousins    & 0.053 & 0.124 & 10 \\
half sibs        & 0.111 & 0.245 & 10 \\
PO / full sibs   & 0.247 & 0.494 & 20 \\
MZ / duplicates  & 0.500 & 1.000 & 10 \\
\bottomrule
\end{tabular}
\end{center}

\noindent
The Pearson correlation between the two is $r=0.988$, and the regression is
$\text{corr.\ resid.} = 2.02\,\hat\phi - 0.004$ --- the factor of two expected if the residual correlation
estimates $2\phi$. Duplicate pairs return exactly $1.000$. Unrelated pairs centre on $-0.003$ with
SD $0.007$.

Two practical consequences. First, if you are running evalAdmix to choose $K$, close relatives will show up
as bad model fit and can be mistaken for a $K$ problem --- so remove relatives first, or read the outliers
as relatedness. Second, in the other direction, an evalAdmix run you already have is a usable relatedness
screen for free.

\textbf{The PCA version} (van Waaij, Li, Garcia-Erill, Albrechtsen \& Wiuf 2023) generalises this to
population structure inferred by \emph{either} PCA or the admixture model, and replaces the original
leave-one-out frequency correction with an analytic bias correction: compare the empirical correlation of
residuals $\hat b$ with the correlation $\hat c$ implied by the model itself, obtained from
$(\mathbf{I}-\mathbf{P})\hat{\mathbf{D}}(\mathbf{I}-\mathbf{P})$ where $\mathbf{P}$ projects onto the column
space of $Q$ (or of the PCs). Under a correct model with $k'=k$, $\hat b - \hat c \to 0$ as the number of
SNPs grows. It is single-pass rather than iterative, needs only genotypes and $Q$ (no ancestral frequency
file), and is available in evalAdmix v1.0 as \texttt{-method corrected}, with an R implementation at
\texttt{github.com/popgenDK/evalPopStructure}. Tellingly, in their 1000 Genomes analysis the elevated
corrected coefficients are explained by five pairs of first-degree relatives --- i.e.\ the statistic
is picking up kinship exactly as above.

\subsection{(C) Avoid allele frequencies altogether}

Use only the joint genotype-sharing pattern within the pair, so that no external frequency --- correct or
incorrect --- can bias the estimate.

\begin{tightitem}
  \item \textbf{KING-robust} (Manichaikul et al.\ 2010) --- kinship from counts of shared heterozygotes
  versus opposing homozygotes. Robust to \emph{discrete} structure and to small samples, and universally used
  as a fast screen. Its documented failure mode is exactly the admixture case: unbiased when both members of
  a pair have similar admixture proportions, systematically negative when they do not.
  \item \textbf{$R_0$, $R_1$, KING-robust kinship} (Waples, Albrechtsen \& Moltke 2019) --- three
  genotype-sharing summary statistics, of which the pairs $(R_1,R_0)$ and $(R_1,\text{KING})$ separate
  parent--offspring, full sibs, half sibs/avuncular/grandparent and unrelated without any allele frequencies
  and without genomic positions. Robust to ascertainment bias; well suited to non-model organisms, ancient
  DNA and reference-free settings. Distinguishes relationship \emph{classes} rather than returning
  unbiased $(k_0,k_1,k_2)$.
\end{tightitem}

\begin{table}[t]
\centering
\small
\caption{Methods for relatedness estimation under admixture / population structure.}
\label{tab:methods}
\begin{tabular}{@{}>{\raggedright\arraybackslash}p{2.3cm} >{\raggedright\arraybackslash}p{1.5cm} >{\raggedright\arraybackslash}p{2.6cm} >{\raggedright\arraybackslash}p{3.1cm} >{\raggedright\arraybackslash}p{4.5cm}@{}}
\toprule
\textbf{Method} & \textbf{Strategy} & \textbf{Input data} & \textbf{Ancestry handled via} & \textbf{Output / notes} \\
\midrule
PLINK \texttt{-{}-genome} & baseline & called genotypes & \emph{nothing} & $(k_0,k_1,k_2)$; biased under admixture \\
NgsRelate & baseline & genotype likelihoods & \emph{nothing} & $(k_0,k_1,k_2)$; single-population, low depth \\
\addlinespace
KING-robust & (C) & called genotypes & pairwise genotype sharing & kinship; fast screen, biased for pairs of \emph{differing} ancestry \\
$R_0/R_1$/KING & (C) & genotypes or GLs & none needed & relationship class; no allele freqs, no positions \\
\addlinespace
REAP & (A) & called genotypes & $\hat q,\hat f$ from ADMIXTURE etc. & $(k_0,k_1,k_2)$, kinship; moment estimator \\
RelateAdmix & (A) & PLINK bed + \texttt{.P}/\texttt{.Q} & $\hat q,\hat f$ from ADMIXTURE etc. & $(k_0,k_1,k_2)$; ML, multithreaded C++ / R \\
InRelate & (A) & multi-allelic genotypes & $\hat q,\hat f$ & all 9 Jacquard states; handles inbreeding \\
SEEKIN & (A) & sparse seq.\ / dosages & ancestry coordinates + ref.\ panel & kinship; works at $\sim$0.15$\times$ depth \\
NGSremix & (A) & genotype likelihoods & $\hat q,\hat f$ (NGSadmix/PCAngsd) & $(k_0,k_1,k_2)$; best for low-depth admixed \\
\addlinespace
PC-AiR & (B) & called genotypes & --- & ancestry PCs unconfounded by relatedness \\
PC-Relate & (B) & called genotypes & PC-AiR PCs & kinship, IBD probs, inbreeding; ``model-free'' \\
PCAngsd & front-end & genotype likelihoods & iterative SVD on dosages $\rightarrow\hat\pi$ & $\hat\pi$, PCA, $\hat q,\hat f$ (NMF), inbreeding; \emph{no} kinship since v1.0 \\
EMU & front-end & called genotypes, much missing & EM-PCA modelling missingness & ancestry PCs unbiased by missingness \\
\addlinespace
evalAdmix & (A) residuals & PLINK bed or Beagle GLs & residual vs.\ $2\hat\pi$ from $\hat q,\hat f$ & corr.\ of residuals $\approx 2\phi$; built as a model-fit test \\
evalAdmix \texttt{-method corrected} & (A)/(B) residuals & PLINK bed & analytic $\hat b-\hat c$ correction, $Q$ or PCs & as above, single-pass, no \texttt{.P} needed \\
\addlinespace
PCAone & front-end & BED/PGEN/BGEN/Beagle/CSV & winSVD; implements EMU \& PCAngsd & fast in-/out-of-core PCA, $\hat\pi$, inbreeding, ancestry-adj. LD; no kinship \\
\bottomrule
\end{tabular}
\end{table}

\section{From kinship to $(k_0,k_1,k_2)$: what PC-Relate actually computes}

One clarification first: \textbf{PC-AiR does not do this}. PC-AiR produces only the ancestry PCs; the IBD
sharing probabilities come from \textbf{PC-Relate}. The description below follows the GENESIS
implementation (\texttt{pcrelate}, \texttt{ibd.probs = TRUE}), where $g_{is}\in\{0,1,2\}$ is the allele count
and $\hat\mu_{is}$ is the individual-specific allele frequency obtained by regressing genotypes on the
PC-AiR PCs (fitted on a training set of mutually unrelated individuals) and halving the prediction.

\paragraph{Kinship.} Form the residual $R_{is}=g_{is}-2\hat\mu_{is}$ and take a normalised cross-product:
\begin{equation*}
\hat\phi_{ij}=\frac{\sum_s R_{is}R_{js}}
{4\sum_s \sqrt{\hat\mu_{is}(1-\hat\mu_{is})\,\hat\mu_{js}(1-\hat\mu_{js})}}.
\end{equation*}

\paragraph{$k_2$ from a dominance coding.} Define the dominance-coded genotype
\begin{equation*}
D_{is}=\begin{cases}\hat\mu_{is} & g_{is}=0\\ 0 & g_{is}=1\\ 1-\hat\mu_{is} & g_{is}=2\end{cases}
\qquad
\tilde D_{is}=D_{is}-\hat\mu_{is}(1-\hat\mu_{is}),
\end{equation*}
which has $E[\tilde D_{is}]=0$ under HWE. Its cross-product isolates the IBD-2 component, because the
dominance covariance between two individuals is proportional to $k_2$ alone:
\begin{equation*}
\hat k_2=\frac{\sum_s \tilde D_{is}\tilde D_{js}}
{\sum_s \hat\mu_{is}(1-\hat\mu_{is})\,\hat\mu_{js}(1-\hat\mu_{js})}.
\end{equation*}

\paragraph{$k_0$ from opposing homozygotes.} If a pair shares \emph{any} allele IBD at a site, they cannot be
opposing homozygotes. So the observed IBS0 count, divided by the count expected if the pair shared nothing,
estimates $k_0$ directly:
\begin{equation*}
\hat k_0=\frac{\sum_s \left[\mathbf{1}(g_{is}{=}2,g_{js}{=}0)+\mathbf{1}(g_{is}{=}0,g_{js}{=}2)\right]}
{\sum_s \left[\hat\mu_{is}^{2}(1-\hat\mu_{js})^{2}+(1-\hat\mu_{is})^{2}\hat\mu_{js}^{2}\right]}.
\end{equation*}

\paragraph{$k_1$ by subtraction.} $\hat k_1 = 1-\hat k_0-\hat k_2$. GENESIS stores \texttt{kin}, \texttt{k0}
and \texttt{k2}; $k_1$ is left as the remainder.

\medskip
\noindent
So the direct answer to ``how does it go from kinship to IBD sharing'' is: \emph{for close pairs, it does
not}. $\hat\phi$, $\hat k_0$ and $\hat k_2$ are three separate estimators computed from the same residuals,
and the identity $\phi=k_1/4+k_2/2$ is not used to derive them. That matters because parent--offspring and
full sibs both have $\phi=1/4$ --- kinship alone cannot tell them apart, and it is the IBS0-based $\hat k_0$
($0$ versus $1/4$) that does.

\paragraph{The corrections, applied only when IBD probabilities are requested.} Three, in order:
\begin{tightitem}
  \item \emph{Inbreeding.} Departure from HWE inflates the dominance covariance, so PC-Relate subtracts the
  product of the two individuals' inbreeding coefficients (the diagonal of the same output):
  $\hat k_2 \leftarrow \hat k_2 - \hat f_i \hat f_j$.
  \item \emph{Residual ancestry and kinship bias in $\hat k_2$ (small-sample correction).} For each PC
  $k\ge 2$, regress the current $\hat k_2$ on the entries of $V_k V_k^{\!\top}$ with a linear and a quadratic
  term, fitting on pairs with $\hat\phi<2^{-11/2}\approx 0.022$ (effectively unrelated), and subtract the
  fitted value. Then regress on $\hat\phi$ itself, fitting on pairs with $\hat k_2<2^{-9/2}\approx 0.044$,
  and subtract that too.
  \item \emph{$k_0$ for more distant pairs.} \emph{Here} the kinship identity finally enters. For pairs below
  the first-degree boundary, $\hat\phi<2^{-5/2}\approx 0.177$, the IBS0 estimate is discarded and replaced by
  \begin{equation*}
  \hat k_0 = 1-4\hat\phi+\hat k_2,
  \end{equation*}
  which is exactly $\phi=k_1/4+k_2/2$ with $k_1=1-k_0-k_2$, rearranged. The reason for the switch is that in
  distant pairs the IBS0 count is dominated by noise and genotyping error, while $\hat\phi$ is still well
  estimated; in first-degree pairs the IBS0 estimator is kept, for the reason above.
\end{tightitem}

\section{Choosing a method}

\begin{tightitem}
  \item \textbf{Array or high-depth genotypes, human cohort, cryptic structure:} PC-AiR + PC-Relate
  (GENESIS). No $K$ to choose, and it is what large consortia use. Screen first with KING-robust.
  \item \textbf{You have a defensible admixture model} (a clear $K$, good reference panels): REAP or
  RelateAdmix. RelateAdmix gives proper ML $(k_0,k_1,k_2)$, which matters if you want to distinguish
  full sibs from parent--offspring rather than just rank kinship.
  \item \textbf{Low-depth sequencing:} do not call genotypes. NGSremix (with NGSadmix or PCAngsd for
  $\hat q,\hat f$), or SEEKIN for very sparse/off-target data.
  \item \textbf{Non-model organism, ancient DNA, no reference frequencies:} $R_0$/$R_1$/KING, which needs
  neither frequencies nor a genetic map.
  \item \textbf{Always:} report which ancestry model was used and at what $K$. Under strategy~(A), the
  relatedness estimates are conditional on $\hat q$ and $\hat f$; a poorly chosen $K$ propagates directly
  into the kinship estimates.
\end{tightitem}

\section{References}
\small
\begin{tightitem}
  \item Manichaikul A, Mychaleckyj JC, Rich SS, Daly K, Sale M, Chen W-M (2010). Robust relationship inference in genome-wide association studies. \emph{Bioinformatics} 26:2867--2873. (KING)
  \item Thornton T, Tang H, Hoffmann TJ, Ochs-Balcom HM, Caan BJ, Risch N (2012). Estimating kinship in admixed populations. \emph{Am J Hum Genet} 91:122--138. (REAP)
  \item Moltke I, Albrechtsen A (2014). RelateAdmix: a software tool for estimating relatedness between admixed individuals. \emph{Bioinformatics} 30:1027--1028.
  \item Conomos MP, Miller MB, Thornton TA (2015). Robust inference of population structure for ancestry prediction and correction of stratification in the presence of relatedness. \emph{Genet Epidemiol} 39:276--293. (PC-AiR)
  \item Conomos MP, Reiner AP, Weir BS, Thornton TA (2016). Model-free estimation of recent genetic relatedness. \emph{Am J Hum Genet} 98:127--148. (PC-Relate)
  \item Dou J, Sun B, Sim X, Hughes JD, Reilly DF, Tai ES, Liu J, Wang C (2017). Estimation of kinship coefficient in structured and admixed populations using sparse sequencing data. \emph{PLoS Genet} 13:e1007021. (SEEKIN)
  \item Meisner J, Albrechtsen A (2018). Inferring population structure and admixture proportions in low-depth NGS data. \emph{Genetics} 210:719--731. (PCAngsd)
  \item Sethuraman A (2018). Estimating genetic relatedness in admixed populations. \emph{G3} 8:3203--3220. (InRelate)
  \item Waples RK, Albrechtsen A, Moltke I (2019). Allele frequency-free inference of close familial relationships from genotypes or low-depth sequencing data. \emph{Mol Ecol} 28:35--48.
  \item Garcia-Erill G, Albrechtsen A (2020). Evaluation of model fit of inferred admixture proportions. \emph{Mol Ecol Resour} 20:936--949. (evalAdmix)
  \item Meisner J, Liu S, Huang M, Albrechtsen A (2021). Large-scale inference of population structure in presence of missingness using PCA. \emph{Bioinformatics} 37:1868--1875. (EMU)
  \item van Waaij J, Li S, Garcia-Erill G, Albrechtsen A, Wiuf C (2023). Evaluation of population structure inferred by principal component analysis or the admixture model. \emph{Genetics} 225:iyad157.
  \item Li Z, Meisner J, Albrechtsen A (2023). Fast and accurate out-of-core PCA framework for large scale biobank data. \emph{Genome Research} 33:1599--1608. (PCAone)
  \item N\o hr AK, Hangh\o j K, Garcia-Erill G, Li Z, Moltke I, Albrechtsen A (2021). NGSremix: a software tool for estimating pairwise relatedness between admixed individuals from next-generation sequencing data. \emph{G3} 11:jkab174.
\end{tightitem}

\end{document}
